% Options for packages loaded elsewhere
\PassOptionsToPackage{unicode}{hyperref}
\PassOptionsToPackage{hyphens}{url}
\PassOptionsToPackage{dvipsnames,svgnames,x11names}{xcolor}
\documentclass[
  11pt,
]{article}
\usepackage{xcolor}
\usepackage[margin=1in]{geometry}
\usepackage{amsmath,amssymb}
\setcounter{secnumdepth}{-\maxdimen} % remove section numbering
\usepackage{iftex}
\ifPDFTeX
  \usepackage[T1]{fontenc}
  \usepackage[utf8]{inputenc}
  \usepackage{textcomp} % provide euro and other symbols
\else % if luatex or xetex
  \usepackage{unicode-math} % this also loads fontspec
  \defaultfontfeatures{Scale=MatchLowercase}
  \defaultfontfeatures[\rmfamily]{Ligatures=TeX,Scale=1}
\fi
\usepackage{lmodern}
\ifPDFTeX\else
  % xetex/luatex font selection
\fi
% Use upquote if available, for straight quotes in verbatim environments
\IfFileExists{upquote.sty}{\usepackage{upquote}}{}
\IfFileExists{microtype.sty}{% use microtype if available
  \usepackage[]{microtype}
  \UseMicrotypeSet[protrusion]{basicmath} % disable protrusion for tt fonts
}{}
\makeatletter
\@ifundefined{KOMAClassName}{% if non-KOMA class
  \IfFileExists{parskip.sty}{%
    \usepackage{parskip}
  }{% else
    \setlength{\parindent}{0pt}
    \setlength{\parskip}{6pt plus 2pt minus 1pt}}
}{% if KOMA class
  \KOMAoptions{parskip=half}}
\makeatother
\usepackage{longtable,booktabs,array}
\newcounter{none} % for unnumbered tables
\usepackage{calc} % for calculating minipage widths
% Correct order of tables after \paragraph or \subparagraph
\usepackage{etoolbox}
\makeatletter
\patchcmd\longtable{\par}{\if@noskipsec\mbox{}\fi\par}{}{}
\makeatother
% Allow footnotes in longtable head/foot
\IfFileExists{footnotehyper.sty}{\usepackage{footnotehyper}}{\usepackage{footnote}}
\makesavenoteenv{longtable}
\setlength{\emergencystretch}{3em} % prevent overfull lines
\providecommand{\tightlist}{%
  \setlength{\itemsep}{0pt}\setlength{\parskip}{0pt}}
\usepackage{booktabs}
\usepackage{microtype}
\usepackage{hyperref}
\setlength{\parskip}{0.4em}
\usepackage{bookmark}
\IfFileExists{xurl.sty}{\usepackage{xurl}}{} % add URL line breaks if available
\urlstyle{same}
\hypersetup{
  pdftitle={OpenCaseLaw: An Open Dataset and Search Platform for Swiss Court Decisions},
  pdfauthor={Jonas Hertner},
  colorlinks=true,
  linkcolor={blue},
  filecolor={Maroon},
  citecolor={Blue},
  urlcolor={blue},
  pdfcreator={LaTeX via pandoc}}

\title{OpenCaseLaw: An Open Dataset and Search Platform for Swiss Court
Decisions}
\author{Jonas Hertner}
\date{March 2026}

\begin{document}
\maketitle

\subsection{Abstract}\label{abstract}

We present OpenCaseLaw, an open corpus and search system for Swiss case
law. The March 20, 2026 snapshot contains 962,724 decisions from 102
federal, cantonal, and quasi-judicial sources, covering all 26 cantons
and the period 1875--2026. The corpus includes 448,461 German (46.6\%),
434,663 French (45.1\%), and 79,600 Italian (8.3\%) decisions.
OpenCaseLaw releases a 34-field Parquet export and reproducible tooling
to build a local search index and reference database. The reference
database contains 8.76 million extracted case-citation references (6.42
million resolved in-corpus, 73.3\% resolution rate) and 11.23 million
decision-to-statute links. We describe the collection, deduplication,
and retrieval pipeline, and release a 100-query multilingual benchmark
with a release-matched baseline report. The code is MIT-licensed; each
record links back to the decision as published by its originating court.

\subsection{1. Introduction}\label{introduction}

Swiss case law is published across a fragmented landscape of federal and
cantonal court websites. Each of the 26 cantons runs its own portal, and
cross-court search does not exist in the open ecosystem. Commercial
systems such as Swisslex and Weblaw aggregate this data but require paid
subscriptions. Existing open resources cover either a single court, a
narrow NLP task, or raw publication access without retrieval tooling.

OpenCaseLaw provides infrastructure rather than a single benchmark. It
combines corpus acquisition, normalization, searchable exports,
reference extraction, and interfaces for both programmatic and
LLM-mediated access. The project contributes:

\begin{enumerate}
\def\labelenumi{\arabic{enumi}.}
\tightlist
\item
  \textbf{A broad open corpus.} The March 20, 2026 snapshot contains
  962,724 decisions from 102 sources---all 26 cantons, federal courts,
  and several regulatory bodies.
\item
  \textbf{Retrieval artifacts.} The release includes Parquet exports
  plus tooling for a local FTS5 search database, a citation and statute
  reference database, REST endpoints, and an MCP server.
\item
  \textbf{Evaluation infrastructure.} A 100-query tagged benchmark and a
  release-matched offline baseline.
\end{enumerate}

The paper reports what is versioned and inspectable in the repository.
Where the implementation distinguishes between data models, search
indexes, and export schemas, we state that explicitly.

\subsection{2. Related Work}\label{related-work}

\textbf{Swiss legal datasets.} Swiss-Judgment-Prediction (Niklaus et
al., 2021) and the Swiss Federal Supreme Court Dataset (Geering and
Merane, 2024) provide Federal Supreme Court resources for downstream
analysis. OpenCaseLaw differs in court coverage: it spans 102 sources
across all cantons and multiple federal and regulatory bodies.

\textbf{Swiss legal benchmarks.} SCALE (Rasiah et al., 2023) benchmarks
citation extraction, court-view generation, and summarization on Swiss
legal text. It is complementary: SCALE focuses on task evaluation,
OpenCaseLaw on corpus and retrieval infrastructure.

\textbf{Broader legal corpora.} MultiLegalPile (Niklaus et al., 2023a)
provides a 689 GB multilingual legal corpus for pretraining; LexGLUE
(Chalkidis et al., 2022) benchmarks legal NLP tasks. Neither provides
Swiss court-wide retrieval infrastructure or citation databases.

\textbf{Open case law infrastructure.} The Caselaw Access Project
(Harvard Law School) provides 6.7 million US court decisions. We are not
aware of a comparable Swiss-wide open corpus with nationwide coverage,
structured metadata, and citation databases.

\textbf{Legal information retrieval.} BM25 and Reciprocal Rank Fusion
remain robust retrieval baselines (Robertson and Zaragoza, 2009; Cormack
et al., 2009). Locke et al.~(2024) survey legal text retrieval.
OpenCaseLaw uses BM25 and RRF as part of a practical multilingual search
pipeline.

\textbf{Re-identification risk.} Pilan et al.~(2024) find that
aggregation and structured metadata increase re-identification risk in
court decisions even when party names are redacted. This applies
directly to corpora like ours.

\subsection{3. Dataset and Processing
Pipeline}\label{dataset-and-processing-pipeline}

\subsubsection{3.1 Snapshot Statistics}\label{snapshot-statistics}

Table 1 reports the frozen release snapshot (March 20, 2026). All
corpus-wide counts in this paper come from this snapshot. Retrieval
metrics in Section 5.3 come from the bundled benchmark artifact.

{\def\LTcaptype{none} % do not increment counter
\begin{longtable}[]{@{}lr@{}}
\toprule\noalign{}
Metric & Value \\
\midrule\noalign{}
\endhead
\bottomrule\noalign{}
\endlastfoot
Snapshot timestamp & 2026-03-20 \\
Decisions & 962,724 \\
Courts and public bodies & 102 \\
Federal sources & 20 \\
Cantonal sources & 82 \\
Federal decisions & 344,141 \\
Cantonal decisions & 618,583 \\
Earliest decision & 1875-01-01 \\
Latest decision & 2026-03-19 \\
German & 448,461 (46.6\%) \\
French & 434,663 (45.1\%) \\
Italian & 79,600 (8.3\%) \\
\end{longtable}
}

The four largest sources are the Federal Supreme Court (174,270), Geneva
(167,003), the Federal Administrative Court (91,613), and Vaud (74,819).

\subsubsection{3.2 Collection}\label{collection}

54 scrapers run nightly, each targeting one court website or publication
portal. Each decision is normalized into a 28-field data model covering
court identity, docket number, dates, language, legal area, regeste,
full text, and source URLs.

\subsubsection{3.3 Deduplication}\label{deduplication}

Deduplication uses a canonical key that normalizes court code, docket
number, and date to collapse formatting variants. Within-court
deduplication keeps the version with the richest content. Cross-court
deduplication operates within hand-maintained overlap groups for cantons
where decisions appear on multiple portals (e.g., Zürich: 17 court
codes).

A BGE leading case and its underlying Federal Supreme Court ruling are
\emph{not} deduplicated---they have different courts, different docket
numbers, and different content scope. Similarly, a cantonal decision and
its federal appeal remain as distinct records.

This is an engineering heuristic, not a curated legal identity model.
The repository does not ship a manual audit of false merges or missed
duplicates.

\subsubsection{3.4 Schemas}\label{schemas}

The repository uses three related schemas:

\begin{enumerate}
\def\labelenumi{\arabic{enumi}.}
\tightlist
\item
  \textbf{Core model} (\texttt{models.py}): 28 fields, used by scrapers.
\item
  \textbf{Search database} (\texttt{db\_schema.py}): 24-column SQLite
  table with FTS5 index.
\item
  \textbf{Parquet export} (\texttt{export\_parquet.py}): 34-field Arrow
  schema with provenance fields.
\end{enumerate}

The project does not have a single monolithic schema; it has a layered
data contract for scraping, retrieval, and export.

\subsection{4. Reference Database}\label{reference-database}

A second SQLite database stores case citations and statute references
extracted from decision text.

\subsubsection{4.1 Extraction}\label{extraction}

Regular expressions extract BGE references
(\texttt{BGE\ 131\ III\ 115}), docket numbers (\texttt{4A\_372/2019}),
and statute provisions (\texttt{Art.\ 41\ OR}). Each reference is
resolved against the corpus using normalized docket matching with
confidence scoring.

\subsubsection{4.2 Scale}\label{scale}

Table 2 summarizes the reference database from the March 20, 2026 build.

{\def\LTcaptype{none} % do not increment counter
\begin{longtable}[]{@{}lr@{}}
\toprule\noalign{}
Metric & Value \\
\midrule\noalign{}
\endhead
\bottomrule\noalign{}
\endlastfoot
Case-citation references & 8.76 million \\
Resolved in-corpus links & 6.42 million \\
Resolution rate & 73.3\% \\
Decision-to-statute links & 11.23 million \\
\end{longtable}
}

These are distinct quantities: 8.76 million refers to case-citation
references; 11.23 million refers to statute links. They should not be
added. The repository does not ship a manual precision/recall study;
Table 2 reports artifact scale, not validated extraction quality.

\subsubsection{4.3 Uses}\label{uses}

The reference database supports citation lookup, leading-case discovery,
appeal-chain tracing, trend analysis, and statute-aware search
enrichment.

\subsection{5. Retrieval and Interfaces}\label{retrieval-and-interfaces}

\subsubsection{5.1 Search Pipeline}\label{search-pipeline}

The search pipeline has five stages:

\begin{enumerate}
\def\labelenumi{\arabic{enumi}.}
\tightlist
\item
  \textbf{Query parsing.} Multiple lexical query variants, legal synonym
  expansion, umlaut normalization, and optional LLM-based structured
  parsing.
\item
  \textbf{Candidate retrieval.} Several FTS5 strategies fused with
  Reciprocal Rank Fusion.
\item
  \textbf{Signal scoring.} Lexical features, metadata, citation counts,
  and statute signals.
\item
  \textbf{Optional reranking.} Confidence-gated LLM reranking for
  ambiguous cases.
\item
  \textbf{Result enrichment.} Court metadata, citation counts, and
  statute references.
\end{enumerate}

\subsubsection{5.2 Distribution}\label{distribution}

The corpus is available as Parquet files, a local SQLite FTS5 index, a
REST API, and an MCP server supporting Claude, ChatGPT, and Gemini. The
MCP tool surface is deployment-dependent (up to 21 tools; remote mode
omits update tools).

\subsubsection{5.3 Evaluation}\label{evaluation}

The repository includes a 100-query benchmark spanning 74 German, 16
French, 7 Italian, and 3 untagged queries across 15 legal domains. Each
query has 1--6 graded relevant decisions (mean 3.04). The judgments were
created by the author; no multi-annotator agreement is reported.

The release bundle includes a benchmark run matched to the
962,724-decision snapshot:

{\def\LTcaptype{none} % do not increment counter
\begin{longtable}[]{@{}lr@{}}
\toprule\noalign{}
Metric & Value \\
\midrule\noalign{}
\endhead
\bottomrule\noalign{}
\endlastfoot
MRR@10 & 0.6042 \\
Recall@10 & 0.5835 \\
nDCG@10 & 0.6062 \\
Hit@1 & 0.52 \\
\end{longtable}
}

This result uses the local search database with the reference graph but
without vector search, statute databases, or LLM-based expansion and
reranking. It is dominated by German queries; concept-match and
statute-oriented retrieval remain the hardest slices.

\subsection{6. Ethics, Legal Basis, and
Limitations}\label{ethics-legal-basis-and-limitations}

\subsubsection{6.1 Legal Basis}\label{legal-basis}

Published Swiss court decisions are excluded from copyright under Art. 5
para. 1 lit. c URG. Federal publication duties are governed by Art. 27
BGG; cantonal duties vary. Code is MIT-licensed; dataset packaging is
CC0-1.0.

OpenCaseLaw preserves decisions as published by the originating courts
and links back to source URLs. It does not perform anonymization.
Cantonal anonymization practices vary.

Large-scale aggregation changes the privacy risk profile. Structured
metadata combined with full text may enable re-identification even when
names are redacted (Pilan et al., 2024). The repository includes a
governance and removal policy.

\subsubsection{6.2 Limitations}\label{limitations}

\begin{itemize}
\tightlist
\item
  \textbf{Coverage is broad, not audited.} Publication depth varies by
  court and era.
\item
  \textbf{Historical quality varies.} Older material may contain OCR
  artifacts.
\item
  \textbf{Reference extraction is rule-based} and misses non-standard
  references.
\item
  \textbf{Identity is operational, not jurisprudential.} Identifiers are
  engineering heuristics, not curated case identities.
\item
  \textbf{No manual evaluation.} Neither citation extraction nor
  deduplication is evaluated against human annotations in this paper.
\item
  \textbf{The benchmark is author-judged.} No inter-annotator agreement
  or held-out test split is provided.
\end{itemize}

\subsection{7. Availability}\label{availability}

All release artifacts are in the GitHub release tagged
\texttt{opencaselaw-paper-2026-03-20}.

{\def\LTcaptype{none} % do not increment counter
\begin{longtable}[]{@{}ll@{}}
\toprule\noalign{}
Resource & Location \\
\midrule\noalign{}
\endhead
\bottomrule\noalign{}
\endlastfoot
Dataset (Parquet) & huggingface.co/datasets/voilaj/swiss-caselaw \\
Source code & github.com/jonashertner/caselaw-repo-1 \\
Release bundle & GitHub release \texttt{opencaselaw-paper-2026-03-20} \\
Governance policy & \texttt{docs/governance-and-removal-policy.md} \\
MCP server & mcp.opencaselaw.ch \\
REST API docs & mcp.opencaselaw.ch/api/docs \\
Dashboard & opencaselaw.ch \\
\end{longtable}
}

\subsection{References}\label{references}

\begin{itemize}
\tightlist
\item
  Caselaw Access Project. Harvard Law School. https://case.law
\item
  Chalkidis, I. et al.~(2022). LexGLUE: A Benchmark Dataset for Legal
  Language Understanding in English. \emph{ACL 2022}.
\item
  Cormack, G., Clarke, C., and Buettcher, S. (2009). Reciprocal Rank
  Fusion outperforms Condorcet and individual Rank Learning Methods.
  \emph{SIGIR 2009}.
\item
  Geering, F. and Merane, J. (2024). Swiss Federal Supreme Court
  Dataset. Zenodo. doi:10.5281/zenodo.11092977.
\item
  Kano, Y. et al.~(2024). COLIEE 2024. \emph{JSAI 2024}.
\item
  Locke, S., Zhai, Z., and Kohlmeier, J. (2024). A Survey on Legal Text
  Retrieval. \emph{ACL 2024}.
\item
  Model Context Protocol. Anthropic, 2024.
  https://modelcontextprotocol.io
\item
  Niklaus, J. et al.~(2021). Swiss-Judgment-Prediction.
  \emph{NLP4PositiveImpact, EMNLP 2021}.
\item
  Niklaus, J. et al.~(2023a). MultiLegalPile: A 689 GB Multilingual
  Legal Corpus. \emph{arXiv:2306.02069}.
\item
  Pilan, I. et al.~(2024). Anonymity at Risk? Re-Identification in Court
  Decisions. \emph{Findings of NAACL 2024}.
\item
  Rasiah, V. et al.~(2023). SCALE: Scaling up Evaluation of Swiss Court
  Rulings. \emph{arXiv:2306.09237}.
\item
  Robertson, S. and Zaragoza, H. (2009). The Probabilistic Relevance
  Framework: BM25 and Beyond. \emph{FTIR}.
\end{itemize}

\end{document}
